% ===== PRÉAMBULE CANONIQUE — à coller VERBATIM en tête de CHAQUE .tex =====
\documentclass[11pt,a4paper]{article}
\usepackage[utf8]{inputenc}\usepackage[T1]{fontenc}\usepackage{lmodern}
\usepackage[french]{babel}
\usepackage{amsmath,amssymb,mathtools}\usepackage{icomma}
\usepackage{siunitx}\sisetup{locale=FR}  % \qty{}{}, \si{}, \ang{}
\usepackage{xcolor}\usepackage{colortbl}
\usepackage{tikz}\usepackage{pgfplots}\pgfplotsset{compat=1.18}
\usepackage[siunitx,europeanresistors]{circuitikz} % circuits (élec.)
\usepackage[most]{tcolorbox}\usepackage{enumitem}\usepackage{array,booktabs}
\usepackage[a4paper,margin=2.2cm]{geometry}
\usetikzlibrary{arrows.meta,calc,angles,quotes,positioning,patterns,%
  backgrounds,shapes.geometric,decorations.pathmorphing,decorations.markings,intersections}
\definecolor{primary}{RGB}{222,49,99}\definecolor{primarydark}{RGB}{150,22,55}
\definecolor{defcol}{RGB}{0,114,178}\definecolor{methcol}{RGB}{52,92,148}
\definecolor{excol}{RGB}{0,135,90}\definecolor{attncol}{RGB}{200,40,55}
\tcbset{boxbase/.style={enhanced,breakable,boxrule=0.9pt,arc=2.5pt,
  left=3mm,right=3mm,top=2.3mm,bottom=2.3mm,fonttitle=\bfseries,coltitle=white}}
% --- boîtes TUTORIEL (noms EXACTS, ne pas renommer) ---
\newtcolorbox{cle}[1][]{boxbase,colback=defcol!6,colframe=defcol,
  title={\textsc{À retenir}\ifx\relax#1\relax\relax\else\textmd{ : \itshape #1}\fi}}
\newtcolorbox{methode}[1][]{boxbase,colback=methcol!7,colframe=methcol,sharp corners,
  title={\textsc{Méthode}\ifx\relax#1\relax\relax\else\textmd{ --- \itshape #1}\fi}}
\newtcolorbox{exemplebox}[1][]{boxbase,colback=excol!5,colframe=excol,
  title={\textsc{Exemple}\ifx\relax#1\relax\relax\else\textmd{ : \itshape #1}\fi}}
\newtcolorbox{piege}[1][]{boxbase,colback=attncol!6,colframe=attncol,
  title={\textsc{Piège}\ifx\relax#1\relax\relax\else\textmd{ : \itshape #1}\fi}}
\newtcolorbox{remarque}[1][]{boxbase,colback=black!4,colframe=black!45,
  title={\textsc{Remarque}\ifx\relax#1\relax\relax\else\textmd{ : \itshape #1}\fi}}
% --- boîtes EXERCICE / CORRIGÉ (noms EXACTS, auto-comptées) ---
\newtcolorbox[auto counter]{exo}[1][]{boxbase,colback=primary!4,colframe=primary,
  title={\textsc{Exercice}~\thetcbcounter\ifx\relax#1\relax\relax\else\textmd{ --- \itshape #1}\fi}}
\newtcolorbox[auto counter]{corrige}[1][]{boxbase,colback=excol!4,colframe=excol,
  title={\textsc{Corrigé}~\thetcbcounter\ifx\relax#1\relax\relax\else\textmd{ --- \itshape #1}\fi}}
\newcommand{\vv}[1]{\vec{#1}}\newcommand{\ds}{\displaystyle}
\newcommand{\R}{\mathbb{R}}\newcommand{\N}{\mathbb{N}}\newcommand{\Z}{\mathbb{Z}}
\begin{document}
% THEME_TITLE: L'œil et ses défauts
\sisetup{per-mode=symbol}

On étudie l'\oe il avec \emph{les mêmes} formules ($C=1/f$, $1/f=1/u+1/v$). Nouveauté : son écran (la
rétine) est à distance \textbf{fixe} --- c'est la \emph{focale} qui s'adapte.

\begin{cle}[le modèle réduit de l'\oe il]
L'\oe il $=$ une \textbf{lentille convergente} (cornée $+$ cristallin, centre optique $O$) qui forme
une image \textbf{réelle et renversée} sur un \textbf{écran} : la \textbf{rétine}. L'\textbf{iris}
règle la lumière entrante (diaphragme). La distance $O$--rétine est \emph{fixe}.
\end{cle}

\begin{center}
\begin{tikzpicture}[scale=0.82,>=Stealth,font=\footnotesize]
  \draw[black!35,dash pattern=on 3pt off 2pt] (-4.6,0)--(4.2,0);
  \draw[line width=1.6pt,<->,black!65] (0,-1.2)--(0,1.2);
  \fill[black] (0,0) circle (1.2pt); \node[above left] at (0,0) {$O$};
  \node[black!70,below,align=center] at (0,-1.25) {cornée\\$+$ cristallin};
  % retine (ecran)
  \fill[attncol!35] (3.2,-1.1) rectangle (3.42,1.1);
  \draw[attncol!80,line width=0.8pt] (3.2,-1.1)--(3.2,1.1);
  \node[attncol!70!black,right,align=left] at (3.45,0.6) {rétine\\(écran)};
  % objet lointain
  \draw[attncol,line width=1.6pt,->] (-4.2,0)--(-4.2,0.9); \node[attncol,left] at (-4.2,0.45) {objet};
  % image renversee sur retine
  \draw[excol,line width=1.6pt,->] (3.2,0)--(3.2,-0.5);
  % rayons paralleles -> convergent sur la retine
  \draw[black!75,line width=0.8pt] (-4.2,0.9)--(0,0.9); \draw[defcol,line width=0.8pt,->] (0,0.9)--(3.2,0);
  \draw[black!75,line width=0.8pt] (-4.2,0)--(0,0);
  \draw[<->,primary,line width=0.6pt] (1.7,1.35)--(3.2,1.35);
  \node[primary,above] at (2.4,1.32) {$f$ fixe};
\end{tikzpicture}
\end{center}

\begin{cle}[accommodation : c'est $f$ qui change, pas la rétine]
Objet \textbf{lointain} : image sur la rétine sans effort. Objet qui \textbf{se rapproche} : le
cristallin \textbf{bombe}, $f$ \emph{diminue}, la vergence $C$ \emph{augmente} pour ramener l'image sur
la rétine. C'est l'\textbf{accommodation}.
\end{cle}

\begin{cle}[les trois défauts et leur correction]
\begin{center}\small
\renewcommand{\arraystretch}{1.25}
\begin{tabular}{>{\bfseries}l l l}
\toprule
Défaut & Image d'un objet\ldots & Correction \\
\midrule
\rowcolor{attncol!8} Myopie & \textbf{en avant} de la rétine (\oe il trop long) & lentille \textbf{divergente} \\
Hypermétropie & \textbf{en arrière} (\oe il trop court) & lentille \textbf{convergente} \\
\rowcolor{defcol!8} Presbytie & \textbf{en arrière} (cristallin durci) & lentille \textbf{convergente} \\
\bottomrule
\end{tabular}
\end{center}
Mnémo : \textbf{myope} $\to$ \textbf{divergente} ; les deux autres $\to$ \textbf{convergente}.
\end{cle}

\begin{methode}[calculer un verre correcteur]
Le verre place l'image \emph{là où l'\oe il voit net} : on applique $1/f=1/u+1/v$ au \emph{verre}, avec
$u$ = position de l'objet et $v$ = point net de l'\oe il (souvent \textbf{virtuel}, $v<0$). Puis
$C=1/f$ (en mètres).
\end{methode}

\begin{exemplebox}[le verre d'un myope]
Myope net au plus loin à \qty{90}{\cm}. Pour voir un objet \textbf{à l'infini}, le verre doit en donner
une image virtuelle à $v=-\qty{0.90}{\m}$ :
$\ \dfrac{1}{f}=0+\dfrac{1}{-0{,}90}=-1{,}11 \Rightarrow f=-\qty{0.90}{\m}$, soit $C=-1{,}11~\delta$
(\textbf{divergent}, comme attendu).
\end{exemplebox}

\begin{piege}[ne pas se tromper de verre]
\begin{itemize}[leftmargin=1.5em,topsep=2pt,itemsep=1pt]
  \item Le \textbf{myope} porte des verres \textbf{divergents} ($C<0$) --- jamais convergents.
  \item La vergence de l'\oe il est \emph{énorme} ($\sim 60~\delta$) : normal, sa focale
        ($\sim\qty{1.6}{\cm}$) est minuscule.
\end{itemize}
\end{piege}

\end{document}
